\chapter{Managing an Exotic Book}\label{ch:dv:managing-an-exotic-book}

A desk sells 100 million of a three-year worst-of \omterm{def:dv:autocallables:autocallable}{autocallable} to its private-bank clients at par. Its model values the notes at 98, so the sale looks
like a profit of 2 million. Valuation control recognises 325\,000 of it. It charges 621\,000 at once for the cost of closing the hedges, and holds back
1\,054\,000 until two things it cannot observe become observable: which model of volatility is right for this payoff, and the correlation between the
two indices. A year later the correlation trades, and 424\,000 comes back. This chapter is about that arithmetic and the book behind it: what an exotic
book owns, the reserves on its valuation, model risk as the desk meets it, stress and concentration, and the \omterm{def:dv:managing-an-exotic-book:dayone}{day-one P\&L} and its release.

\section{What an exotic book owns}

An exotic book is a large short position in \omterm{def:dv:the-structured-products-business:product}{structured products} (chapters 18 and 19) and the hedges bought against it, vanillas, variance and
the underlying. The net \omterm{def:dv:greeks-and-the-hedging-pnl:greeks}{Greeks} are small; what remains are the exposures nobody sells. For \omterm{def:dv:autocallables:autocallable}{autocallables} sold to investors, that is long volatility at long
expiries, long skew and short correlation on worst-ofs, all concentrated where the notes' barriers sit.

Accounting sorts those exposures by the inputs their valuation needs. IFRS 13 ranks inputs in three levels: quoted prices of identical instruments in
active markets (Level 1); other inputs observable directly or indirectly, for substantially the full term of the instrument (Level 2); unobservable inputs
(Level 3). A three-year correlation between two indices has no market: the worst-of note is a Level 3 instrument, and the rules on its day-one profit follow
from that.

\begin{example}[The note]\label{ex:dv:managing-an-exotic-book:note}
Worst-of Phoenix \omterm{def:dv:autocallables:autocallable}{autocallable} on two indices (volatilities 20\% and 25\%, correlation 0.5, rates 3\%, dividends 2\%), three annual observations,
autocall at 100\%, a coupon of 7.95\% a year paid at or above 70\% with memory, capital protected at maturity unless the worse index is below 60\%.
At that coupon the model values the note at 98.00 per 100. From the desk's side, short the note, the book is long 0.533 and 0.693 of \omterm{def:dv:greeks-and-the-hedging-pnl:greeks}{vega} per volatility
point in the two indices (per 100 of notional) and short 0.143 and 0.279 of delta per 1\%, which the desk hedges by buying the indices.
\end{example}

\section{Reserves and parameter bid--offer}

The booked value is a mid. What the firm could exit at is worse, for three reasons, each a reserve.

\begin{definition}[Valuation reserve]\label{def:dv:managing-an-exotic-book:reserve}
A \emph{valuation reserve}\index{valuation reserve} is an amount by which a firm lowers the booked value of a position below its model mid, for a
specific source of uncertainty in what the position could be exited at; reserves are computed position by position, kept in the books, and reviewed at
every valuation date.
\end{definition}

\begin{definition}[Bid--offer reserve]\label{def:dv:managing-an-exotic-book:bidoffer}
A \emph{bid--offer reserve}\index{bid--offer reserve} is the cost of closing a book's net risk at the market's bid or offer rather than at mid: the net
exposure in each risk bucket times half the bid--offer spread of the instruments that would close it.
\end{definition}

\begin{definition}[Model reserve]\label{def:dv:managing-an-exotic-book:model}
A \emph{model reserve}\index{model reserve} covers the uncertainty from the choice of model: the booked model's value less a prudent point of the values
the position takes under a set of alternative models that market participants use.
\end{definition}

\begin{definition}[Parameter bid--offer]\label{def:dv:managing-an-exotic-book:param}
The \emph{parameter bid--offer}\index{parameter bid--offer} of a position is the reserve for an unobservable input: the booked value less a prudent point of
the values over the input's plausible range.
\end{definition}

The EU's prudent-valuation standard fixes the prudent point for regulatory capital: where a range of plausible values exists, the point where the firm is
90\% confident it could exit at that value or better. The build uses the same rule for all three reserves.

\begin{example}[Three reserves on the note]\label{ex:dv:managing-an-exotic-book:reserves}
\emph{Bid--offer.} Closing the \omterm{def:dv:greeks-and-the-hedging-pnl:greeks}{vega} at half a point of spread and the delta at 2 basis points costs 0.266 and 0.346 for the two \omterm{def:dv:greeks-and-the-hedging-pnl:greeks}{vegas} and less than 0.01 for the deltas: 0.621 in all
per 100. \emph{Model.} Four models value the note at 98.00 (flat at-the-money volatilities, booked), 94.33 (flat at the knock-in strike's volatility, three points
higher), 98.62 (flat at the trigger's, half a point lower) and 97.98 (uncertain volatility, $\pm20\%$). The prudent point is the highest of these for a short
position, 98.62: a reserve of 0.621. \emph{Correlation.} Over a plausible range of 0.35 to 0.65 the note is worth 97.64 to 98.50, and the 90\% point is 98.43: a
reserve of 0.433 (\cref{fig:dv:managing-an-exotic-book:rho,fig:dv:managing-an-exotic-book:models}).
\end{example}

\begin{omfigure}
\begin{tikzpicture}
\begin{axis}[width=10.5cm,height=5cm,xlabel={correlation between the two indices},ylabel={note value (per 100)},
  tick label style={font=\small},label style={font=\small},
  xmin=0.18,xmax=0.82,ymin=97.2,ymax=99.2,grid=major,grid style={gray!25},
  table/col sep=comma]
\fill[omProp!15] (axis cs:0.35,97.2) rectangle (axis cs:0.65,99.2);
\fill[omThm!30] (axis cs:0.47,97.2) rectangle (axis cs:0.53,99.2);
\addplot[omThm,very thick,mark=*,mark size=1.4pt] table[x=rho,y=value]{figdata/derivatives/27-managing-an-exotic-book/rho.csv};
\addplot[black,dashed] coordinates {(0.18,98.433) (0.82,98.433)};
\node[font=\footnotesize,anchor=south west] at (axis cs:0.19,98.44) {90\% point at inception};
\end{axis}
\end{tikzpicture}

\omcaption{Value of the three-year worst-of note against the correlation of its two indices: the plausible range at inception (light, 0.35 to 0.65) and once the
correlation is observable (dark, 0.47 to 0.53). The firm is short the note, so higher correlation is the unfavourable side. Data: the
tutorial.}\label{fig:dv:managing-an-exotic-book:rho}
\end{omfigure}

\begin{omfigure}
\begin{tikzpicture}
\begin{axis}[width=10.5cm,height=5cm,xbar,bar width=10pt,xlabel={note value (per 100)},
  tick label style={font=\small},label style={font=\small},
  xmin=93,xmax=99.5,ymin=-0.6,ymax=3.6,ytick={0,1,2,3},
  yticklabels={flat at-the-money (booked),flat at the knock-in strike,flat at the trigger,uncertain volatility},
  grid=major,grid style={gray!25},table/col sep=comma]
\addplot[omDef,fill=omDef!40] table[x=value,y=i]{figdata/derivatives/27-managing-an-exotic-book/models.csv};
\draw[black,thick,dashed] (axis cs:98.621,-0.6) -- (axis cs:98.621,3.6);
\end{axis}
\end{tikzpicture}

\omcaption{The note under four volatility models. For the short position the prudent point (dashed) is the highest value, half a volatility point of \omterm{def:dv:greeks-and-the-hedging-pnl:greeks}{vega} above the
booked model. Data: the tutorial.}\label{fig:dv:managing-an-exotic-book:models}
\end{omfigure}

The reserves are not added naively for capital. The EU standard's first aggregation method counts half of each position's market-price, close-out and model-risk
adjustments, for diversification across positions: 0.838 here.

\section{Model risk on the desk}

Book 6 treats model risk as a governance process. On the desk it has three faces. The first is the choice among models that all fit the vanillas and disagree on
the exotic, the \omterm{def:dv:managing-an-exotic-book:model}{model reserve}'s range: here almost four points of note value between the flat models at two strikes. The second is a model used outside its
domain, such as a flat-volatility pricer for a knock-in whose value lives in the skew. The third is a calibration that moves the value without a market move
(chapter 24's \omterm{def:dv:fourier-pricing-and-calibration-engineering:recalib}{recalibration P\&L}).

\begin{dated}{2026-09}{Model-risk guidance in the United States}\label{dat:dv:managing-an-exotic-book:mrm}
On 17 April 2026 the Federal Reserve, the OCC and the FDIC replaced the 2011 model-risk guidance (SR 11-7) with revised guidance (Federal Reserve letter SR 26-2).
It defines a model as ``a complex quantitative method, system, or approach that applies statistical, economic, or financial theories to process input data into
quantitative estimates'', and leaves generative and agentic AI models out of its scope.
\end{dated}

A desk lives with model risk by pricing every exotic under at least two models, reserving the difference, and hedging what the models agree on. When the models
disagree about the hedge (the delta of a worst-of under two correlation assumptions, the \omterm{def:dv:greeks-and-the-hedging-pnl:greeks}{vega} of a barrier under local and stochastic volatility), the desk hedges
the part they agree on and limits the rest.

\section{Stress and concentration}

The \omterm{def:dv:greeks-and-the-hedging-pnl:greeks}{Greeks} describe small moves. A book of short \omterm{def:dv:autocallables:autocallable}{autocallables} has its risk in large ones: the knock-in barriers cluster, and the hedges that were right at inception
turn over as the indices approach them.

\begin{example}[A stress grid]\label{ex:dv:managing-an-exotic-book:stress}
The firm's P\&L on the delta-hedged short note, in millions for 100 million of notional, under joint moves of both indices and of both volatilities:
\begin{center}\small
\begin{tabular}{lrrrrr}
\toprule
& $-30\%$ & $-20\%$ & $-10\%$ & $0$ & $+10\%$\\
\midrule
volatilities $-5$ points & 7.02 & $-2.27$ & $-6.31$ & $-5.89$ & $-2.85$\\
unchanged & 13.09 & 5.63 & 1.31 & 0.00 & 0.96\\
volatilities $+5$ points & 18.21 & 12.15 & 8.08 & 6.08 & 5.83\\
\bottomrule
\end{tabular}
\end{center}
The firm is long the investors' knock-in put, so a crash with rising volatility is a gain. The loss is a moderate fall with falling volatility: $-6.31$ million at
$-10\%$ and five points less volatility (\cref{fig:dv:managing-an-exotic-book:stress}).
\end{example}

\begin{omfigure}
\begin{tikzpicture}
\begin{axis}[width=10.5cm,height=5cm,xlabel={move of both indices (\%)},ylabel={P\&L (millions)},
  tick label style={font=\small},label style={font=\small},
  xmin=-32,xmax=12,ymin=-8,ymax=20,grid=major,grid style={gray!25},
  legend columns=3,legend cell align=left,legend style={font=\footnotesize,at={(0.5,-0.3)},anchor=north,draw=none,fill=none,/tikz/every even column/.append style={column sep=8pt}},
  table/col sep=comma]
\addplot[omProp,very thick,mark=*,mark size=1.4pt] table[x=shock,y=v-5]{figdata/derivatives/27-managing-an-exotic-book/stress.csv};
\addplot[black,thick,mark=square*,mark size=1.4pt] table[x=shock,y=v+0]{figdata/derivatives/27-managing-an-exotic-book/stress.csv};
\addplot[omThm,very thick,mark=triangle*,mark size=1.6pt] table[x=shock,y=v+5]{figdata/derivatives/27-managing-an-exotic-book/stress.csv};
\legend{volatilities $-5$ points,unchanged,$+5$ points}
\end{axis}
\end{tikzpicture}

\omcaption{The delta-hedged short worst-of note under joint moves of the indices and their volatilities (100 million of notional). The book's worst case is a moderate
fall with falling volatility, not a crash. Data: the tutorial.}\label{fig:dv:managing-an-exotic-book:stress}
\end{omfigure}

Concentration is the other face of size. The book's \omterm{def:dv:greeks-and-the-hedging-pnl:greeks}{vega} in the second index, 693\,000 per point, is a large position against the long-dated volatility that trades.
At an assumed 100\,000 per point a day and a participation of a tenth, closing it would take 69 days. The prudent-valuation standard asks for a concentration
adjustment when the prudent exit period exceeds ten days. When many banks sell the same notes, the same exposure sits on every book, and the exit is crowded. Losses
on such books are public when they are large. Natixis reported a non-recurring impact of \euro259 million on Asian equity derivatives in the fourth quarter of 2018,
after identifying a hedging strategy in Asia as deficient.

\section{Day-one P\&L and its release}

\begin{definition}[Day-one P\&L]\label{def:dv:managing-an-exotic-book:dayone}
The \emph{day-one P\&L}\index{day-one P\&L} of a trade is the difference, at inception, between its transaction price and its fair value; under IFRS 9 it is
recognised at once only if the fair value is evidenced by a Level 1 price or by a valuation using only observable data, and is otherwise deferred and recognised
later only as a factor that market participants would price, including time, changes.
\end{definition}

The rule does not say how to split a margin whose valuation uses some observable and some unobservable inputs. The build's policy is one reading. The \omterm{def:dv:managing-an-exotic-book:bidoffer}{bid--offer
reserve}, on observable inputs, is part of fair value from the first day. The reserves on unobservable inputs, model and correlation, are deferred. What remains of the
margin after both is recognised. A stricter policy defers the whole margin until the last unobservable input can be observed.

\begin{example}[The deferred profit]\label{ex:dv:managing-an-exotic-book:dayone}
Per 100 million of notes sold at par: margin 2.000 million; \omterm{def:dv:managing-an-exotic-book:bidoffer}{bid--offer reserve} 0.621; deferred 1.054 (model 0.621, correlation 0.433); recognised at inception
0.325. After one year, with both indices at 95\% and the note alive, the correlation trades within 0.47--0.53 and the note has two years left. The correlation
reserve falls to 0.071 and the \omterm{def:dv:managing-an-exotic-book:model}{model reserve} to 0.559: 0.424 million is released, 0.362 of it from the correlation. After two years, with the indices still at 95\%,
another 0.336 comes back, and the last 0.295 at maturity (\cref{fig:dv:managing-an-exotic-book:release}).
\end{example}

\begin{omfigure}
\begin{tikzpicture}
\begin{axis}[width=10.5cm,height=5cm,ybar,bar width=14pt,xlabel={year},ylabel={millions},
  tick label style={font=\small},label style={font=\small},
  xmin=-0.6,xmax=3.6,ymin=0,ymax=1.2,xtick={0,1,2,3},grid=major,grid style={gray!25},
  legend columns=2,legend cell align=left,legend style={font=\footnotesize,at={(0.5,-0.3)},anchor=north,draw=none,fill=none,/tikz/every even column/.append style={column sep=8pt}},
  table/col sep=comma]
\addplot[omDef,fill=omDef!40] table[x=year,y=deferred]{figdata/derivatives/27-managing-an-exotic-book/release.csv};
\addplot[omThm,fill=omThm!40] table[x=year,y=released]{figdata/derivatives/27-managing-an-exotic-book/release.csv};
\legend{deferred balance,released in the year}
\end{axis}
\end{tikzpicture}

\omcaption{The deferred part of the note's day-one margin: 1.054 million held back at inception, released as the correlation becomes observable and the note shortens
(indices assumed at 95\% at each anniversary). Data: the tutorial.}\label{fig:dv:managing-an-exotic-book:release}
\end{omfigure}

Deferral is not a loss. It moves profit from the year the note was sold to the years its uncertainty is resolved. It changes behaviour all the same: a desk paid on
recognised P\&L is less keen on trades whose margin lives in unobservable inputs.

\section{Tutorial: reserving an autocallable}\label{tut:dv:managing-an-exotic-book:tut}

\begin{tutorial}
\textbf{Goal.} Price a worst-of \omterm{def:dv:autocallables:autocallable}{autocallable}, set its coupon for a 2\% margin, compute its bid--offer, model and correlation reserves, split its \omterm{def:dv:managing-an-exotic-book:dayone}{day-one P\&L} and
schedule its release, and run a stress grid. \textbf{End state:} the four figures, the stress table and the numbers of the weekend problem.
\begin{tutsteps}
\item \textbf{The prudent point} and the model and parameter reserves:
\omcode{code/firm/reserves/firm_reserves.py}{22}{47}{Prudent point, model reserve and parameter bid--offer.}\label{lst:dv:managing-an-exotic-book:prudent}
\item \textbf{The day-one split} and the release schedule:
\omcode{code/firm/reserves/firm_reserves.py}{56}{81}{Day-one P\&L and its release.}\label{lst:dv:managing-an-exotic-book:dayone}
\item \textbf{Run} \texttt{dv\_reserves.fair\_coupon()}, \texttt{reserves\_at(0)}, \texttt{day\_one\_study()}, \texttt{rho\_curve()}, \texttt{stress()} and
\texttt{fig\_reserves.py}.
\end{tutsteps}
\textbf{What to change next.} Add chapter 18's local-volatility pricer to the model set; value the release with the indices at 80\% after a year; set the correlation
range from the dispersion of \omterm{def:dv:multi-asset-options:implied}{implied correlations} of chapter 17.
\end{tutorial}

\section{Build: the reserve calculator}\label{bld:dv:managing-an-exotic-book:reserves}

\begin{build}
\bfield{Purpose} The miniature firm's reserve calculator: bid--offer by risk bucket, \omterm{def:dv:managing-an-exotic-book:model}{model reserves} from a model set, \omterm{def:dv:managing-an-exotic-book:param}{parameter bid--offer} on unobservable inputs,
aggregation, the day-one split and its release, stress grids and concentration.
\bfield{Interface} \texttt{bid\_offer\_reserve(exposures, half\_spreads)}; \texttt{prudent\_point(values, confidence, weights)}; \texttt{model\_reserve(values\_by\_model,
booked)}; \texttt{parameter\_reserve(value\_of, mid, lo, hi)}; \texttt{aggregate(reserves, diversification)}; \texttt{day\_one(price, booked\_liability, observable,
unobservable)} $\to$ \texttt{DayOne}; \texttt{release\_schedule}; \texttt{stress\_grid(value\_of, spot\_shocks, vol\_shocks, hedge\_delta)}; \texttt{concentration\_days}.
\bfield{Rules} Values are the book's value to the firm and reserves are positive; the prudent point is the 90\% point of the plausible range; unobservable reserves are
deferred, observable ones are part of fair value; every reserve is recomputed at every valuation date.
\bfield{Acceptance tests} \texttt{code/firm/reserves/tests/}: the prudent point of a uniform range; a \omterm{def:dv:managing-an-exotic-book:model}{model reserve} of zero when the booked model is already the prudent one;
the parameter reserve of a linear value on a uniform range; bid--offer, aggregation and day-one arithmetic; a release schedule that sums to the deferral; a stress grid on a
linear book.
\bfield{Stretch} Reserves across a book with netting of offsetting exposures; a concentration adjustment from the exit period; a model set from chapters 9, 10 and 18's
pricers through chapter 28's library.
\end{build}

\begin{omsources}
\item Commission Regulation (EU) No 1255/2012 (IFRS 13 \emph{Fair Value Measurement}), paragraphs 76--86.
\item Commission Regulation (EU) 2016/2067 (IFRS 9 \emph{Financial Instruments}), paragraph B5.1.2A.
\item Commission Delegated Regulation (EU) 2016/101 (prudent valuation), Articles 9, 11 and 14 and Annex.
\item Board of Governors of the Federal Reserve System, SR 26-2, ``Revised guidance on model risk management'' (17 April 2026); OCC Bulletin 2026-13.
\item Natixis, 2018 fourth-quarter and annual results, press release, 12 February 2019.
\end{omsources}

\section{Exercises}

\begin{exercise}[$\star$]\label{exo:dv:managing-an-exotic-book:1}
Why is the worst-of note a Level 3 instrument, and when could it become Level 2?
\end{exercise}

\begin{exercise}[$\star$]\label{exo:dv:managing-an-exotic-book:2}
Why is the prudent point the highest of the model values for the firm, which is short the note?
\end{exercise}

\begin{exercise}[$\star$]\label{exo:dv:managing-an-exotic-book:3}
Check that the day-one split adds up, and compute it under a policy that defers the whole margin net of bid--offer.
\end{exercise}

\begin{exercise}[$\star\star$]\label{exo:dv:managing-an-exotic-book:4}
The \omterm{def:dv:managing-an-exotic-book:model}{model reserve} (0.621) equals the \omterm{def:dv:greeks-and-the-hedging-pnl:greeks}{vega} bid--offer (0.621) to three decimals. Is this double counting?
\end{exercise}

\begin{exercise}[$\star\star$]\label{exo:dv:managing-an-exotic-book:5}
Why does higher correlation raise the value of a worst-of note?
\end{exercise}

\begin{exercise}[$\star\star$]\label{exo:dv:managing-an-exotic-book:6}
Explain why the stress grid's worst cell is a moderate fall with falling volatility.
\end{exercise}

\begin{exercise}[$\star\star\star$]\label{exo:dv:managing-an-exotic-book:7}
\emph{Coding.} Compute the correlation reserve with the plausible range 0.2--0.8 instead of 0.35--0.65. What does the \omterm{def:dv:managing-an-exotic-book:dayone}{day-one P\&L} become?
\end{exercise}

\begin{exercise}[$\star\star\star$]\label{exo:dv:managing-an-exotic-book:8}
\emph{Find the flaw.} ``Our reserves are conservative: we add all of them for every position, with no aggregation benefit, so capital is never understated.''
\end{exercise}

\section{Problem: The Deferred Profit}

\begin{problem}[{Weekend problem --- how much of the margin is profit}]\label{pb:dv:managing-an-exotic-book:1}
The desk of the opening sells 100 million of the three-year worst-of \omterm{def:dv:autocallables:autocallable}{autocallable} at par; its model values the notes at 98.

\medskip\noindent\textbf{Part I --- The note.}
\begin{enumerate}
\item Give the coupon that makes the note worth 98, and say what the firm is long and short once it has sold the note.
\item Which inputs of the valuation are observable, and which are not?
\item At what level of the fair-value hierarchy is the note, and what does IFRS 9 say about its day-one gain?
\item Give the \omterm{def:dv:greeks-and-the-hedging-pnl:greeks}{vegas} and deltas of the firm's position.
\item What does the firm hedge, and what can it not hedge?
\end{enumerate}

\medskip\noindent\textbf{Part II --- Reserves.}
\begin{enumerate}[resume]
\item Compute the \omterm{def:dv:managing-an-exotic-book:bidoffer}{bid--offer reserve}.
\item Give the four models' values and the \omterm{def:dv:managing-an-exotic-book:model}{model reserve}.
\item Give the note's range of values over the correlation's plausible range and the correlation reserve.
\item What is the aggregated total under the EU standard's first method?
\item How long would it take to close the second index's \omterm{def:dv:greeks-and-the-hedging-pnl:greeks}{vega}, and what follows?
\end{enumerate}

\medskip\noindent\textbf{Part III --- Day one and after.}
\begin{enumerate}[resume]
\item Split the margin into charged, deferred and recognised.
\item What changes after one year, and what is released?
\item How much of the release comes from the correlation?
\item What would a fall of the indices to 80\% in the first year do to the release?
\item How would the stress grid's worst cell change the desk's hedging?
\end{enumerate}

\medskip\noindent\textbf{Part IV --- Judgement.}
\begin{enumerate}[resume]
\item Is deferring the margin conservative or just slow?
\item How does deferral change the desk's incentives?
\item Who should own the model set and the plausible ranges?
\item State the \emph{named result}: the \omterm{def:dv:managing-an-exotic-book:dayone}{day-one P\&L} recognised at inception on a three-year worst-of \omterm{def:dv:autocallables:autocallable}{autocallable}, and the amount released after one year as the
correlation input becomes observable.
\item In one sentence: what is a reserve for?
\end{enumerate}
\end{problem}

\section{Interview questions}

\begin{interviewq}[$\star$ \iqroles{bank, risk}]\label{iq:dv:managing-an-exotic-book:1}
What is \omterm{def:dv:managing-an-exotic-book:dayone}{day-one P\&L}, and when can a bank recognise it?
\end{interviewq}

\begin{interviewq}[$\star\star$ \iqroles{risk}]\label{iq:dv:managing-an-exotic-book:2}
How would you compute a \omterm{def:dv:managing-an-exotic-book:model}{model reserve} for an exotic?
\end{interviewq}

\begin{interviewq}[$\star\star$ \iqroles{bank, trader}]\label{iq:dv:managing-an-exotic-book:3}
What exposures does a bank that issues \omterm{def:dv:autocallables:autocallable}{autocallables} keep, and why are they hard to hedge?
\end{interviewq}

\begin{interviewq}[$\star\star$ \iqroles{risk}]\label{iq:dv:managing-an-exotic-book:4}
Design a stress test for a book of short \omterm{def:dv:autocallables:autocallable}{autocallables}. Which scenario do you fear most?
\end{interviewq}

\begin{interviewq}[$\star\star$ \iqroles{developer, risk}]\label{iq:dv:managing-an-exotic-book:5}
What must a reserve calculation system store so that an auditor can reproduce a reserve a year later?
\end{interviewq}

\begin{interviewq}[$\star\star\star$ \iqroles{bank}]\label{iq:dv:managing-an-exotic-book:6}
An unobservable correlation becomes observable. What happens to the book's P\&L, and why?
\end{interviewq}
